% !TEX root = ../main.tex
\section{Patrones de diseño}\label{sec:patrones}

Esta sección registra los problemas del juego en los que se evaluó aplicar un patrón, la decisión tomada en cada uno y, para los patrones aplicados, su ficha. Un patrón se aplicó solo cuando existía una fuerza concreta que lo justificara; en dos de los problemas analizados la decisión fue no aplicar ninguno. Los patrones de estructura y creacionales se identifican aquí como candidatos y se aplican en la Entrega~3 \citar{gamma2003patrones}.

\subsection{Tabla de decisión}

{\tablacompacta
\begin{longtable}{@{}C{0.6cm}L{3.3cm}L{2.7cm}L{2.2cm}L{5.7cm}@{}}
  \caption{Problemas, patrones candidatos y decisión tomada}\label{tab:decision}\\
  \toprule
  \cab{N.\textordmasculine} & \cab{Problema y fuerza} & \cab{Candidatos} & \cab{Decisión} & \cab{Por qué se descartaron los otros} \\ \midrule \endfirsthead
  \toprule
  \cab{N.\textordmasculine} & \cab{Problema y fuerza} & \cab{Candidatos} & \cab{Decisión} & \cab{Por qué se descartaron los otros} \\ \midrule \endhead
  \bottomrule \endlastfoot
  D1 & Cada nivel de CPU elige su jugada con un criterio distinto, y se espera agregar niveles sin tocar la partida. & Strategy; Template Method; subclases de \texttt{JugadorCpu} (diseño 1.0); condicional por nivel. & \textbf{Strategy} (ficha en la sección~\ref{sec:ficha-strategy}). & Template Method supone un esqueleto común con pasos variables, pero cada nivel reemplaza la decisión completa. Las subclases atan el criterio al tipo de jugador y obligan a crear una subclase por nivel. Un condicional concentra en una clase criterios que no se relacionan. \\ \filasep
  D2 & Las acciones admitidas cambian según el estado de la partida: preparada, salida, en curso o terminada. & State; condicionales sobre \texttt{EstadoPartida} (diseño 1.0); Chain of Responsibility. & \textbf{State} (ficha en la sección~\ref{sec:ficha-state}). & Los condicionales obligan a consultar el mismo atributo en cada una de las siete operaciones de \texttt{Partida} que dependen del estado, y un estado nuevo obliga a revisarlas todas. Chain of Responsibility encadena validaciones, pero aquí la regla depende de cuál estado está activo, no de una secuencia de comprobaciones. \\ \filasep
  D3 & Las jugadas deben poder deshacerse (RN-10) y la partida debe poder guardarse y reconstruirse (RN-11). & Command; Memento; Prototype. & \textbf{Command} (ficha en la sección~\ref{sec:ficha-command}). & Memento y Prototype guardan instantáneas del estado: resuelven deshacer, pero el archivo guardado expondría el estado interno de la partida y no conservaría qué se jugó, que la vista necesita mostrar. Command usa la misma secuencia de jugadas para deshacer, para mostrar las últimas jugadas y para reconstruir la partida al cargar. \\ \filasep
  D4 & Validar una jugada antes de ejecutarla: turno, estado, ficha en la mano, coincidencia con el extremo y robo permitido. & Chain of Responsibility; validación en el estado y en la jugada. & \textbf{Sin patrón adicional.} & Son cuatro comprobaciones fijas y baratas, sin variación que aislar; una cadena de manejadores agregaría clases sin una fuerza que las justifique. El estado decide qué jugadas admite y cada jugada verifica sus propias condiciones con \texttt{EsValida}. \\ \filasep
  D5 & Informar a la interfaz las jugadas que hicieron las CPU. & Observer; consulta de la vista después de cada operación. & \textbf{Sin patrón adicional.} & Hay un único interesado, la interfaz o el simulador, que recibe la vista como respuesta de cada operación, y \texttt{VistaPartida.UltimasJugadas} ya contiene las jugadas recientes. Observer agregaría suscripción y orden de notificación sin un segundo suscriptor. Se reconsiderará si se agrega un marcador o un registro de eventos. \\ \filasep
  D6 & La interfaz necesita un único punto de entrada que oculte la estructura del dominio. & Facade; acceso directo a \texttt{Partida}. & \textbf{Facade} en \texttt{Controlador\allowbreak Juego}; se documenta en la Entrega~3. & El acceso directo acopla la interfaz a las reglas y obliga a repetir la coordinación de los turnos de las CPU en la consola y en la simulación. \\ \filasep
  D7 & Integrar el reloj del sistema y el almacenamiento en archivos JSON sin que el dominio dependa de ellos. & Adapter; uso directo de \texttt{DateTime.\allowbreak Now} y del sistema de archivos. & \textbf{Adapter} en \texttt{Reloj\allowbreak Sistema} y \texttt{Repositorio\allowbreak Partidas\allowbreak Json}; se documenta en la Entrega~3. & El uso directo impide la simulación con reloj controlado y ata las reglas al medio de almacenamiento. \\ \filasep
  D8 & Crear los jugadores y su estrategia a partir de la configuración. & Factory Method; Abstract Factory; condicional en \texttt{Partida.\allowbreak Crear\allowbreak Jugadores} (actual). & \textbf{Pendiente} para la Entrega~3. & Hoy un condicional elige la estrategia según el nivel. Se evaluará en la unidad de patrones creacionales, junto con la construcción de la partida. \\
\end{longtable}}

% ---------------------------------------------------------------------
\clearpage
\subsection{Ficha del patrón Strategy}\label{sec:ficha-strategy}

{\small\setstretch{1.05}\renewcommand{\arraystretch}{1.2}
\begin{longtable}{@{}L{3.1cm}L{11.4cm}@{}}
  \caption{Ficha del patrón Strategy}\label{tab:ficha-strategy}\\
  \toprule
  \cab{Campo} & \cab{Contenido} \\ \midrule \endfirsthead
  \toprule
  \cab{Campo} & \cab{Contenido} \\ \midrule \endhead
  \bottomrule \endlastfoot
  Patrón y unidad & Strategy. Patrón de comportamiento, Unidad~3. \\ \filasep
  Problema y fuerzas & Los niveles de CPU solo difieren en cómo eligen entre las opciones válidas: al azar, por más puntos o conservando fichas jugables. En la versión 1.0 cada nivel era una subclase de \texttt{JugadorCpu}: el criterio quedaba fijado al crear el jugador y cada nivel nuevo exigía otra subclase. Fuerzas: agregar criterios sin modificar \texttt{Partida}, probarlos por separado y decidir sobre \texttt{VistaPartida}, sin ver las manos ajenas. \\ \filasep
  Clases participantes & \emph{Strategy}: \texttt{IEstrategiaCpu}. \emph{ConcreteStrategy}: \texttt{EstrategiaPrincipiante}, \texttt{EstrategiaIntermedia}, \texttt{EstrategiaAvanzada}. \emph{Context}: \texttt{JugadorCpu}, que delega \texttt{DecidirJugada} y \texttt{Nivel} en su estrategia. \emph{Cliente}: \texttt{Partida.CrearJugadores}. \\ \filasep
  Estructura & Figura~\ref{fig:patron-strategy}. \\ \filasep
  Consecuencias & \textbf{Mejora:} una sola clase de jugador automático; cada criterio es una clase pequeña, sin estado, que se prueba con una vista construida en la prueba; el nivel ya no se duplica, porque lo da la estrategia. \textbf{Costo:} una interfaz y una indirección más por decisión, y una clase más en total. Elegir la estrategia según el nivel sigue siendo un condicional en \texttt{CrearJugadores}, que se trata en la Entrega~3 (decisión D8). \\ \filasep
  Alternativa descartada & Template Method: no hay un esqueleto común con pasos que varíen, porque cada nivel reemplaza la decisión completa, y conservaría la herencia que ata el criterio al jugador. \\ \filasep
  Prueba de que sirve & Un nivel nuevo es una clase que implementa \texttt{IEstrategiaCpu} más un caso en \texttt{CrearJugadores}, sin tocar \texttt{JugadorCpu}, \texttt{ResolverTurnosCpu} ni la interfaz; antes exigía una subclase de jugador. En 1000 partidas simuladas con semilla 42 los criterios se distinguen: Avanzado gana 296, Intermedio 252, el humano simulado 229 y Principiante 215 (8 empates). El costo de la indirección se medirá en la Entrega~3. \\
\end{longtable}}

\diagrama{patron_strategy}{Strategy: estructura antes y después de aplicarlo}{fig:patron-strategy}

% ---------------------------------------------------------------------
\clearpage
\subsection{Ficha del patrón State}\label{sec:ficha-state}

{\small\setstretch{1.05}\renewcommand{\arraystretch}{1.2}
\begin{longtable}{@{}L{3.1cm}L{11.4cm}@{}}
  \caption{Ficha del patrón State}\label{tab:ficha-state}\\
  \toprule
  \cab{Campo} & \cab{Contenido} \\ \midrule \endfirsthead
  \toprule
  \cab{Campo} & \cab{Contenido} \\ \midrule \endhead
  \bottomrule \endlastfoot
  Patrón y unidad & State. Patrón de comportamiento, Unidad~3. \\ \filasep
  Problema y fuerzas & Lo que la partida admite depende de su estado: preparada, solo iniciar; salida, solo la ficha de salida y sin tiempo; en curso, colocar, robar, pasar y deshacer con control de tiempo; terminada, nada. En la versión 1.0 el estado era un atributo \texttt{EstadoPartida} que cada operación debía consultar con condicionales, como ya lo hacían \texttt{Iniciar} y \texttt{OpcionesDe}. Fuerzas: las reglas de un estado quedaban repartidas en siete operaciones, un estado nuevo obligaba a revisarlas todas y las transiciones quedaban implícitas. \\ \filasep
  Clases participantes & \emph{State}: \texttt{EstadoJuego}. \emph{ConcreteState}: \texttt{EstadoPreparada}, \texttt{EstadoSalida}, \texttt{EstadoEnCurso}, \texttt{EstadoTerminada}. \emph{Context}: \texttt{Partida}, que delega en su estado las siete operaciones y ofrece \texttt{CambiarEstado}. \texttt{EstadoPartida} queda como valor visible en la vista. \\ \filasep
  Estructura & Figuras~\ref{fig:patron-state} (antes y después) y~\ref{fig:estados-partida} (transiciones). \\ \filasep
  Consecuencias & \textbf{Mejora:} las reglas de cada estado están juntas, las transiciones son explícitas y \texttt{Partida} perdió los condicionales; por omisión un estado no admite nada, así que uno nuevo solo habilita lo que declara. \textbf{Costo:} cinco clases más; los estados usan operaciones de paquete de \texttt{Partida}, lo que los acopla al contexto, y cada acción queda repartida entre la partida y el estado. \\ \filasep
  Alternativa descartada & Mantener los condicionales de la versión 1.0 y Chain of Responsibility (decisión D2). Tampoco se conservó Cerrada como estado: duraba un instante y no tenía reglas propias; el cierre queda en \texttt{Resultado.Forma}. \\ \filasep
  Prueba de que sirve & La regla de que el tiempo no corre en la salida (RN-04) se expresó solo en los estados, sin tocar \texttt{Partida}: \texttt{EstadoEnCurso} controla el tiempo y \texttt{EstadoSalida} hereda el comportamiento por omisión. Deshacer la salida es una transición a \texttt{EstadoSalida}. Las pruebas de \texttt{EstadoPruebas} verifican lo que admite cada estado. \\
\end{longtable}}

\diagrama{patron_state}{State: estructura antes y después de aplicarlo}{fig:patron-state}

\diagrama{estados_partida}{Diagrama de estados de la partida}{fig:estados-partida}

% ---------------------------------------------------------------------
\clearpage
\subsection{Ficha del patrón Command}\label{sec:ficha-command}

{\small\setstretch{1.05}\renewcommand{\arraystretch}{1.2}
\begin{longtable}{@{}L{3.1cm}L{11.4cm}@{}}
  \caption{Ficha del patrón Command}\label{tab:ficha-command}\\
  \toprule
  \cab{Campo} & \cab{Contenido} \\ \midrule \endfirsthead
  \toprule
  \cab{Campo} & \cab{Contenido} \\ \midrule \endhead
  \bottomrule \endlastfoot
  Patrón y unidad & Command. Patrón de comportamiento, Unidad~3. \\ \filasep
  Problema y fuerzas & Deshacer (RN-10) y guardar y cargar (RN-11) exigen conservar las jugadas en orden, revertirlas y repetirlas. En la versión 1.0 \texttt{Jugada} era un dato que \texttt{Partida.Aplicar} y \texttt{Revertir} debían interpretar con un condicional. Fuerzas: ejecutar y revertir deben ser simétricos; robar exige recordar la ficha; cada tipo nuevo tocaba dos operaciones de \texttt{Partida}; el archivo no debe exponer el estado interno. \\ \filasep
  Clases participantes & \emph{Command}: \texttt{Jugada}. \emph{ConcreteCommand}: \texttt{JugadaColocarFicha}, \texttt{JugadaRobar}, \texttt{JugadaPasar}, \texttt{JugadaTiempoAgotado}. \emph{Receiver}: \texttt{Mesa}, \texttt{Mano}, \texttt{FichasParaRobar}. \emph{Invoker}: \texttt{Partida} (\texttt{Aplicar}, \texttt{RevertirHastaUltimaDe}, \texttt{Restaurar}). \texttt{HistorialJugadas} conserva los comandos ejecutados. \\ \filasep
  Estructura & Figura~\ref{fig:patron-command}. \\ \filasep
  Consecuencias & \textbf{Mejora:} cada jugada se ejecuta y se deshace sobre sus receptores; la misma secuencia sirve para deshacer, mostrar y guardar; un tipo nuevo no toca \texttt{Partida}; la interfaz lee las jugadas sin poder ejecutarlas. \textbf{Costo:} cinco clases reemplazan un tipo de valor; \texttt{JugadaRobar} no es inmutable; el JSON exige registrar los tipos en el repositorio; cargar vuelve a procesar todas las jugadas, acción que se medirá. \\ \filasep
  Alternativa descartada & Memento: una instantánea de mesa, manos y fichas antes de cada jugada del humano resolvía deshacer, pero no guardar sin exponer el estado interno ni conservar qué se jugó. Prototype, clonar la partida cada turno, tiene el mismo problema y más memoria (decisión D3). \\ \filasep
  Prueba de que sirve & Un tipo de jugada nuevo es una subclase de \texttt{Jugada} y su registro en el repositorio, sin tocar deshacer, guardar ni \texttt{Restaurar}. \texttt{JugadaPruebas} verifica que deshacer restaura el estado, incluida la salida, y que cargar reconstruye la partida, también tras deshacer un robo. La simulación deshace, guarda y carga en cada una de sus 1000 partidas. \\
\end{longtable}}

\diagrama{patron_command}{Command: estructura antes y después de aplicarlo}{fig:patron-command}
